% Canal Box — Documentation technique du projet
% Compiler : pdflatex CanalBox_Documentation.tex (×2 pour la TOC)
% ou : python manage.py generate_project_tex

\documentclass[11pt,a4paper]{article}

\usepackage[T1]{fontenc}
\usepackage[utf8]{inputenc}
\usepackage[french]{babel}
\usepackage{geometry}
\usepackage{xcolor}
\usepackage{tikz}
\usepackage{tcolorbox}
\tcbuselibrary{skins,breakable}
\usepackage{booktabs}
\usepackage{array}
\usepackage{tabularx}
\usepackage{enumitem}
\usepackage{hyperref}
\usepackage{fancyhdr}
\usepackage{graphicx}
\usepackage{fontawesome5}
\usepackage{listings}
\usepackage{titlesec}
\usepackage{parskip}

\geometry{margin=2.2cm, top=2.8cm, bottom=2.5cm}

% ── Palette Canal Box ──────────────────────────────────────────────
\definecolor{cbPink}{HTML}{ff2d78}
\definecolor{cbPinkDark}{HTML}{e0155f}
\definecolor{cbPinkSoft}{HTML}{fff0f6}
\definecolor{cbInk}{HTML}{0f0f12}
\definecolor{cbSlate}{HTML}{4b5565}
\definecolor{cbMist}{HTML}{98a2b3}
\definecolor{cbPos}{HTML}{12b76a}
\definecolor{cbNeg}{HTML}{f04438}
\definecolor{cbNeu}{HTML}{f79009}
\definecolor{cbGlass}{HTML}{f8f9fc}
\definecolor{cbDark}{HTML}{1a1a1f}
\definecolor{cbPurple}{HTML}{6941c6}

% ── Hyperliens ─────────────────────────────────────────────────────
\hypersetup{
  colorlinks=true,
  linkcolor=cbPink,
  urlcolor=cbPinkDark,
  citecolor=cbSlate,
  pdftitle={Canal Box — Documentation},
  pdfauthor={Équipe Deep Learning LSTM},
}

% ── En-têtes / pieds de page ───────────────────────────────────────
\pagestyle{fancy}
\fancyhf{}
\fancyhead[L]{\small\color{cbSlate}Canal Box — Documentation technique}
\fancyhead[R]{\small\color{cbPink}\thepage}
\renewcommand{\headrulewidth}{2pt}
\renewcommand{\headrule}{\hbox to\headwidth{\color{cbPink}\leaders\hrule height \headrulewidth\hfill}}

% ── Titres colorés ─────────────────────────────────────────────────
\titleformat{\section}{\Large\bfseries\color{cbPink}}{}{0em}{}[\vspace{-0.3em}\color{cbPink}\titlerule]
\titleformat{\subsection}{\large\bfseries\color{cbInk}}{}{0em}{}

% ── Boîtes colorées ────────────────────────────────────────────────
\newtcolorbox{archbox}[2][]{
  enhanced, breakable,
  colback=#2!8, colframe=#2,
  fonttitle=\bfseries\color{white},
  coltitle=white, title=#1,
  attach boxed title to top left={yshift=-2mm, xshift=4mm},
  boxed title style={colback=#2, sharp corners},
  arc=4pt, boxrule=0.8pt,
  left=8pt, right=8pt, top=10pt, bottom=8pt,
}

\newtcolorbox{infobox}[1][]{
  enhanced, breakable,
  colback=cbPinkSoft, colframe=cbPink,
  fonttitle=\bfseries\color{cbPinkDark},
  title={\faInfoCircle\ #1},
  arc=4pt, boxrule=0.6pt,
  left=8pt, right=8pt,
}

\newtcolorbox{warnbox}[1][]{
  enhanced, breakable,
  colback=cbNeg!8, colframe=cbNeg,
  fonttitle=\bfseries\color{cbNeg},
  title={\faShieldAlt\ #1},
  arc=4pt, boxrule=0.6pt,
}

% ── Listings (arborescence) ─────────────────────────────────────────
\lstset{
  basicstyle=\ttfamily\small\color{cbSlate},
  backgroundcolor=\color{cbGlass},
  frame=single, rulecolor=\color{cbMist},
  frameround=tttt, framesep=6pt,
  breaklines=true, showstringspaces=false,
  columns=fullflexible, keepspaces=true,
}

% ── Badges sentiment ───────────────────────────────────────────────
\newcommand{\badgePos}{\tikz[baseline=(X.base)]{\node[fill=cbPos, text=white, rounded corners=2pt, inner xsep=5pt, inner ysep=2pt, font=\small\bfseries] (X) {Positif};}}
\newcommand{\badgeNeg}{\tikz[baseline=(X.base)]{\node[fill=cbNeg, text=white, rounded corners=2pt, inner xsep=5pt, inner ysep=2pt, font=\small\bfseries] (X) {Négatif};}}
\newcommand{\badgeNeu}{\tikz[baseline=(X.base)]{\node[fill=cbNeu, text=white, rounded corners=2pt, inner xsep=5pt, inner ysep=2pt, font=\small\bfseries] (X) {Neutre};}}

% ── Page de titre ──────────────────────────────────────────────────
\newcommand{\makecover}{
  \begin{titlepage}
    \begin{tikzpicture}[remember picture, overlay]
      \fill[cbDark] (current page.south west) rectangle (current page.north east);
      \fill[cbPink, opacity=0.35] ([xshift=14cm, yshift=24cm]current page.south west) circle (3.5cm);
      \fill[cbPink, opacity=0.15] ([xshift=2cm, yshift=4cm]current page.south west) circle (2.2cm);
    \end{tikzpicture}
    \centering
    \vspace*{2.5cm}
    {\tikz{\node[fill=cbPink, rounded corners=8pt, minimum width=1.8cm, minimum height=1.8cm, font=\Huge\bfseries\color{white}] {C};}}\\[1.2cm]
    {\Huge\bfseries\color{white} Canal Box}\\[0.4cm]
    {\Large\color{cbPinkSoft} Plateforme d'avis clients · Analyse LSTM}\\[1.5cm]
    {\large\color{cbMist} Documentation technique du projet}\\[0.3cm]
    {\color{cbMist} Examen Deep Learning — Cahier des charges v1.1}\\[0.2cm]
    {\color{cbMist} Juillet 2026}\\[2cm]
    \begin{tikzpicture}
      \foreach \x/\lbl/\col in {0/Django 5.2/cbPink, 2.2/SQLite/cbSlate, 4.4/LSTM/cbPos, 6.6/Liquid Glass/cbNeu} {
        \node[fill=\col, text=white, rounded corners=3pt, inner xsep=8pt, inner ysep=4pt, font=\small\bfseries] at (\x,0) {\lbl};
      }
    \end{tikzpicture}
    \vfill
    {\small\color{cbMist} Équipe Deep Learning LSTM · Usage académique}
  \end{titlepage}
}

% ═══════════════════════════════════════════════════════════════════
\begin{document}
\makecover
\tableofcontents
\newpage

% ── 1. Présentation ────────────────────────────────────────────────
\section{Présentation du projet}

Canal Box est une plateforme simulée d'avis clients pour un opérateur de services
numériques (internet fibre, décodeur TV, assistance client, facturation).
Chaque avis soumis est automatiquement analysé par un réseau \textbf{LSTM} qui prédit
la polarité du texte : \badgePos\ \badgeNeg\ \badgeNeu.

\begin{table}[h]
  \centering
  \renewcommand{\arraystretch}{1.4}
  \begin{tabularx}{\textwidth}{>{\bfseries\color{cbPink}}l X}
    \toprule
    \rowcolor{cbPink}\color{white}Objectif & \color{white}\bfseries Description \\
    \midrule
    Collecte & Permettre aux clients de publier et consulter des avis \\
    Analyse IA & Classifier automatiquement chaque avis via LSTM \\
    Administration & Tableau de bord, modération, logs, exports CSV/PDF \\
    Confidentialité IA & Résultats LSTM visibles \textbf{uniquement} par les administrateurs \\
    \bottomrule
  \end{tabularx}
  \caption{Objectifs principaux}
\end{table}

% ── 2. Architecture ────────────────────────────────────────────────
\section{Architecture technique}

L'application repose sur une architecture \textbf{monolithique Django} : templates,
logique métier, inférence LSTM et persistance SQLite dans un seul projet.

\begin{center}
\begin{tikzpicture}[
  layer/.style={rectangle, rounded corners=6pt, minimum width=12cm, minimum height=1.1cm, align=center, font=\small},
  arr/.style={-{Stealth[length=3mm]}, thick, cbPink}
]
  \node[layer, fill=cbPinkSoft, draw=cbPink, text=cbInk] (pres) at (0,4)
    {\textbf{Présentation}\quad|\quad Templates Django · Tailwind · GSAP · Liquid Glass};
  \node[layer, fill=cbGlass, draw=cbSlate, text=cbInk] (met) at (0,2.5)
    {\textbf{Métier}\quad|\quad views.py · review\_service · stats\_service · presenters};
  \node[layer, fill=cbPos!12, draw=cbPos, text=cbInk] (ia) at (0,1)
    {\textbf{Intelligence artificielle}\quad|\quad lstm\_service → lstm\_stub / lstm\_keras (.h5)};
  \node[layer, fill=blue!6, draw=blue!50, text=cbInk] (data) at (0,-0.5)
    {\textbf{Données}\quad|\quad SQLite · User · Category · Review · SentimentResult · InferenceLog};
  \draw[arr] (pres) -- (met);
  \draw[arr] (met) -- (ia);
  \draw[arr] (ia) -- (data);
\end{tikzpicture}
\end{center}

\begin{table}[h]
  \centering
  \renewcommand{\arraystretch}{1.35}
  \begin{tabularx}{\textwidth}{>{\bfseries}l >{\ttfamily\small}l X}
    \toprule
    \rowcolor{cbInk}\color{white}Composant & \color{white}\bfseries Technologie & \color{white}\bfseries Rôle \\
    \midrule
    Backend & Django 5.2 + ORM & Vues, auth, services métier \\
    Base de données & SQLite & Persistance locale des avis et résultats IA \\
    Frontend & Django Templates + Tailwind & Interface Liquid Glass responsive \\
    Animations & GSAP + ScrollTrigger & Transitions et micro-interactions \\
    IA (actif) & lstm\_stub.py & Heuristiques FR en attendant le modèle \\
    IA (futur) & lstm\_keras.py + .h5 & Modèle TensorFlow/Keras entraîné \\
    \bottomrule
  \end{tabularx}
  \caption{Stack technologique}
\end{table}

% ── 3. Flux ────────────────────────────────────────────────────────
\section{Flux de données — soumission d'un avis}

À chaque soumission, le pipeline suivant s'exécute de manière transparente pour le client :

\begin{enumerate}[leftmargin=*, itemsep=6pt]
  \item[\tikz{\node[circle, fill=cbPink, text=white, font=\bfseries\small, minimum size=6mm]{1};}]
    \textbf{Client} — Rédige un avis (titre, texte, note 1–5, catégorie)
  \item[\tikz{\node[circle, fill=cbSlate, text=white, font=\bfseries\small, minimum size=6mm]{2};}]
    \textbf{Backend} — Enregistre en base (\texttt{status: pending})
  \item[\tikz{\node[circle, fill=cbPos, text=white, font=\bfseries\small, minimum size=6mm]{3};}]
    \textbf{LSTM} — Analyse le texte → sentiment + score de confiance
  \item[\tikz{\node[circle, fill=cbSlate, text=white, font=\bfseries\small, minimum size=6mm]{4};}]
    \textbf{Backend} — Met à jour (\texttt{status: analyzed}) + log d'inférence
  \item[\tikz{\node[circle, fill=cbPinkSoft, draw=cbPink, text=cbPink, font=\bfseries\small, minimum size=6mm]{5};}]
    \textbf{Client} — Voit l'avis publié \textbf{sans} étiquette de sentiment
  \item[\tikz{\node[circle, fill=cbDark, text=white, font=\bfseries\small, minimum size=6mm]{6};}]
    \textbf{Administrateur} — Consulte sentiment, confiance, stats et modération
\end{enumerate}

\newpage

% ── 4. Modèle de données ───────────────────────────────────────────
\section{Modèle de données}

\begin{table}[h]
  \centering
  \renewcommand{\arraystretch}{1.4}
  \begin{tabularx}{\textwidth}{>{\ttfamily\small}l >{\ttfamily\small}l X}
    \toprule
    \rowcolor{cbPinkDark}\color{white}\bfseries Table & \color{white}\bfseries Champs clés & \color{white}\bfseries Description \\
    \midrule
    users & email, role, plan, status & Comptes client et administrateur \\
    categories & name, slug, description, icon & Services Canal Box évaluables \\
    reviews & user\_id, title, content, rating, status & Avis clients soumis \\
    sentiment\_results & sentiment, confidence, prob\_* & Résultat LSTM par avis \\
    inference\_logs & status, latency\_ms, model\_version & Journal des inférences \\
    \bottomrule
  \end{tabularx}
  \caption{Schéma relationnel (SQLite / ORM Django)}
\end{table}

\begin{archbox}[User]{cbPurple}
  \texttt{role} : \texttt{client} | \texttt{admin} \quad
  \texttt{status} : \texttt{actif} | \texttt{suspendu}
\end{archbox}

\begin{archbox}[Review → SentimentResult (1:1)]{cbPos}
  \texttt{status} : \texttt{pending} → \texttt{analyzed} \quad
  \texttt{flagged} : signalé si confiance $<$ seuil minimal
\end{archbox}

% ── 5. Structure ───────────────────────────────────────────────────
\section{Structure des dossiers}

\begin{lstlisting}
canalbox/              --> settings.py, urls.py, wsgi.py
core/
  models.py            --> User, Category, Review, SentimentResult
  views.py             --> 30+ vues (public, client, admin)
  presenters.py        --> ORM -> dicts pour templates
  services/
    review_service.py  --> Creation avis + inference LSTM
    stats_service.py   --> Statistiques dashboard admin
  ml/
    lstm_service.py    --> Interface d'inference
    lstm_stub.py       --> Stub heuristique (actif)
    lstm_keras.py      --> Branche modele reel (a venir)
  management/commands/
    seed_canalbox.py   --> Jeu de donnees initial
templates/             --> public/, client/, adminpanel/, partials/
static/                --> canalbox.css, canalbox.js
ml_models/             --> lstm_sentiment.h5 (a venir)
docs/                  --> Documentation PDF et LaTeX
\end{lstlisting}

% ── 6. Pages ───────────────────────────────────────────────────────
\section{Pages de l'application}

\subsection{Espace public}

\begin{table}[h]
  \centering
  \small
  \renewcommand{\arraystretch}{1.3}
  \begin{tabularx}{\textwidth}{>{\ttfamily}l X >{\bfseries\color{cbPurple}}l}
    \toprule
    \rowcolor{cbPurple}\color{white}\bfseries URL & \color{white}\bfseries Page & \color{white}\bfseries Accès \\
    \midrule
    / & Landing — présentation Canal Box & Tous \\
    /connexion/ & Connexion e-mail / mot de passe & Visiteurs \\
    /inscription/ & Création de compte client & Visiteurs \\
    /mot-de-passe-oublie/ & Réinitialisation mot de passe & Visiteurs \\
    \bottomrule
  \end{tabularx}
\end{table}

\subsection{Espace client}

\begin{table}[h]
  \centering
  \small
  \renewcommand{\arraystretch}{1.3}
  \begin{tabularx}{\textwidth}{>{\ttfamily}l X >{\bfseries\color{cbPink}}l}
    \toprule
    \rowcolor{cbPink}\color{white}\bfseries URL & \color{white}\bfseries Page & \color{white}\bfseries Accès \\
    \midrule
    /espace/ & Tableau de bord personnel & Client \\
    /espace/avis/ & Liste, recherche, filtres, pagination & Client \\
    /espace/avis/nouveau/ & Formulaire de soumission d'avis & Client \\
    /espace/avis/\textless id\textgreater/ & Détail d'un avis (sans IA) & Client \\
    /espace/profil/ & Profil et statistiques & Client \\
    /espace/profil/modifier/ & Modification du profil & Client \\
    \bottomrule
  \end{tabularx}
\end{table}

\subsection{Console administrateur}

\begin{table}[h]
  \centering
  \small
  \renewcommand{\arraystretch}{1.3}
  \begin{tabularx}{\textwidth}{>{\ttfamily}l X >{\bfseries\color{cbInk}}l}
    \toprule
    \rowcolor{cbInk}\color{white}\bfseries URL & \color{white}\bfseries Page & \color{white}\bfseries Accès \\
    \midrule
    /administration/ & Dashboard KPI + graphiques sentiments & Admin \\
    /administration/avis/ & Liste avec sentiment \& confiance LSTM & Admin \\
    /administration/avis/\textless id\textgreater/ & Détail + distribution softmax & Admin \\
    /administration/moderation/ & Avis signalés (confiance faible) & Admin \\
    /administration/utilisateurs/ & Gestion comptes (suspendre) & Admin \\
    /administration/categories/ & CRUD catégories de services & Admin \\
    /administration/logs/ & Journal inférences (latence, erreurs) & Admin \\
    /administration/exports/csv/ & Export CSV des avis analysés & Admin \\
    /administration/exports/pdf/ & Export rapport analytique & Admin \\
    \bottomrule
  \end{tabularx}
\end{table}

\newpage

% ── 7. Sécurité ────────────────────────────────────────────────────
\section{Règles de sécurité \& rôles}

\begin{warnbox}[Règle fondamentale — cahier des charges]
  Les données IA (\texttt{sentiment}, \texttt{confidence}, logs, statistiques LSTM) ne sont
  \textbf{jamais} transmises aux templates client. Seuls les administrateurs y accèdent.
\end{warnbox}

\begin{table}[h]
  \centering
  \renewcommand{\arraystretch}{1.35}
  \begin{tabularx}{\textwidth}{>{\bfseries\color{cbNeg}}l X}
    \toprule
    \rowcolor{cbNeg}\color{white}Règle & \color{white}\bfseries Implémentation \\
    \midrule
    IA invisible côté client & \texttt{review\_to\_dict(include\_ia=False)} \\
    Accès admin protégé & Décorateur \texttt{admin\_required} + \texttt{role == admin} \\
    Comptes suspendus & Vérification \texttt{status} avant chaque action \\
    Seuil de confiance & \texttt{MIN\_CONFIDENCE\_THRESHOLD} → flagged automatique \\
    Validation entrées & Longueur min titre/contenu, note 1–5, e-mail unique \\
    \bottomrule
  \end{tabularx}
\end{table}

% ── 8. LSTM ────────────────────────────────────────────────────────
\section{Intégration du modèle LSTM}

Le stub actuel (\texttt{lstm\_stub.py}) utilise des heuristiques sur mots-clés français.
Pour brancher le modèle entraîné par le pôle IA :

\begin{infobox}[Étapes d'intégration]
\begin{enumerate}[leftmargin=*, nosep]
  \item Exporter le modèle depuis Jupyter (format \texttt{.h5} ou SavedModel)
  \item Placer le fichier dans \texttt{ml\_models/lstm\_sentiment.h5}
  \item Implémenter \texttt{predict()} dans \texttt{core/ml/lstm\_keras.py}
  \item Changer \texttt{LSTM\_BACKEND = 'keras'} dans \texttt{canalbox/settings.py}
\end{enumerate}
\end{infobox}

\textbf{Contrainte cahier des charges :} temps d'inférence $<$ 3 secondes par avis.

% ── 9. Installation ────────────────────────────────────────────────
\section{Installation \& commandes}

\begin{table}[h]
  \centering
  \renewcommand{\arraystretch}{1.35}
  \begin{tabularx}{\textwidth}{>{\ttfamily\small}l X}
    \toprule
    \rowcolor{cbPos}\color{white}\bfseries Commande & \color{white}\bfseries Description \\
    \midrule
    pip install -r requirements.txt & Installer les dépendances \\
    python manage.py migrate & Créer les tables SQLite \\
    python manage.py seed\_canalbox & Insérer données de démonstration \\
    python manage.py runserver & Lancer le serveur de développement \\
    python manage.py generate\_project\_pdf & Régénérer le PDF \\
    pdflatex CanalBox\_Documentation.tex & Compiler ce document LaTeX \\
    \bottomrule
  \end{tabularx}
\end{table}

\begin{infobox}[Comptes de démonstration]
  \textbf{Client} : \texttt{yanice.client@canalbox.cd} \\
  \textbf{Admin} : \texttt{david.admin@canalbox.cd} \\
  Mot de passe : \texttt{demo1234}
\end{infobox}

\vfill
\begin{center}
  \color{cbMist}\small
  \rule{0.3\textwidth}{2pt}\\[0.5em]
  © 2026 Canal Box — Équipe Deep Learning LSTM\\
  Document confidentiel — Usage académique
\end{center}

\end{document}
